\chapter{Mesh, Materials, and Load Assembly}

On the current parallel PMG path, preprocessing still serves the same purpose
as in MATLAB: turn the mesh and material table into the objects needed by the
constitutive builder and continuation solver. The important difference is that
the default path no longer starts by building one global \code{B} matrix and
one global elastic matrix.

\section{Mainline Preprocessing Sequence}

For the committed 3D benchmark, \petscmod{cli/run\_3D\_hetero\_SSR\_capture.py}
performs the following sequence:
\begin{enumerate}
\item load the Gmsh mesh through \code{load\_mesh\_from\_file(...)},
\item reorder nodes with \code{block\_metis},
\item compute \code{n\_q} from \code{quadrature\_volume\_3d(elem\_type)},
\item expand material rows over integration points through
\code{heterogenous\_materials(...)},
\item choose the lightweight MPI elastic path,
\item assemble owned elastic rows through
\code{assemble\_owned\_elastic\_rows\_for\_comm(...)},
\item convert those rows into a PETSc AIJ matrix \code{K\_elast},
\item gather the local right-hand-side pieces into the global body-force field
\code{f\_V}.
\end{enumerate}

So the current default path reaches the same physical objects as MATLAB, but it
reaches them through rank-owned assembly rather than a global sparse build.

\section{Object Mapping on the Mainline Path}

\begin{longtable}{@{}p{0.22\textwidth}p{0.34\textwidth}p{0.38\textwidth}@{}}
\toprule
Object & MATLAB meaning & Mainline PETSc meaning \\
\midrule
\endhead
\code{coord} & nodal coordinates & reordered coordinates aligned with PETSc ownership \\
\code{elem} & element connectivity & same connectivity after node permutation \\
\code{surf} & boundary faces & same boundary meaning after node permutation \\
\code{q\_mask} & free displacement mask & same mask, reused later by solver and free-DOF extraction \\
\code{material\_identifier} & material row id per element & same element-wise id after reorder \\
\code{c0}, \code{phi}, \code{psi}, \code{shear}, \code{bulk}, \code{lame}, \code{gamma} & integration-point material arrays & same arrays built by \code{heterogenous\_materials(...)} \\
\code{K\_elast} & elastic stiffness matrix & PETSc AIJ matrix assembled from owned rows \\
\code{f\_V} & body-force vector & global field formed by gathering rank-local body-force pieces \\
\bottomrule
\end{longtable}

\section{Node Reordering Is Part of The Mainline}

For the current benchmark, node reordering is not an optional PETSc tuning
detail. It is part of the default path because the runner sets
\code{node\_ordering = block\_metis}. That means:
\begin{itemize}
\item the numbering is chosen to align with MPI ownership,
\item later owned-row matrix assembly uses that numbering directly,
\item the constitutive and tangent chapters inherit this reordered ownership map.
\end{itemize}

MATLAB did not need an explicit public reordering stage because its global
matrix lived in one address space.

\section{Material Expansion Is Still A Direct Translation}

\petscmod{mesh/materials.py::heterogenous\_materials} is still the closest
translation of \matlabpkg{+ASSEMBLY/heterogenous\_materials.m}. It still:
\begin{itemize}
\item interprets material ids element-wise,
\item converts \code{phi} and \code{psi} to radians,
\item derives \code{shear}, \code{bulk}, and \code{lame},
\item repeats material rows over the element quadrature rule,
\item builds the integration-point weight term \code{gamma}.
\end{itemize}

For the non-seepage mainline run, the saturation mask is effectively fixed, so
the mechanical body-force path is determined directly from the material table
without any preceding seepage solve.

\section{What The Mainline Path Does Not Build}

The most important difference from the earlier document version is this:
\begin{itemize}
\item the current default PMG path does \textbf{not} use the global
\code{assemble\_strain\_operator(...) + build\_elastic\_stiffness\_matrix(...)}
route as its normal startup path,
\item instead it builds owned elastic rows immediately and lets the owned
tangent pattern carry the detailed strain geometry for later Newton updates.
\end{itemize}

The global reference assembly is still available, but it belongs in the annex
because it is not the default parallel PMG route.

\section{Where To Edit}

\begin{itemize}
\item \petscmod{mesh/loader.py} for mesh loading,
\item \petscmod{mesh/reorder.py} for ownership-oriented numbering,
\item \petscmod{mesh/materials.py} for integration-point material expansion,
\item \petscmod{fem/distributed\_elastic.py} for the owned elastic precursor,
\item \petscmod{cli/run\_3D\_hetero\_SSR\_capture.py} for the exact mainline
assembly order.
\end{itemize}

\section{Where The Other Preprocessing Paths Went}

The global reference assembly, other node orderings, and seepage-oriented mesh
loaders are moved to Appendix~\ref{app:mesh}.
